\section{Methodology}
\label{sec:method}

In this section, we present the mathematical formulation and structural design of \textbf{SE2RoadNet}.
We first define the planar curve representation and formalize the requirement for $\mathrm{SE}(2)$ rigid-motion invariance.
We then detail the 7-channel coordinate-free feature set, the relative-arclength attention bias mechanism, and the physics-informed regularizer.

\subsection{Problem Formulation and Invariance Guarantee}

Let a test scenario road trajectory be represented as a continuous planar curve $\mathbf{r}(s) = (x(s), y(s))^\top \in \mathbb{R}^2$, parameterized by its arclength $s \in [0, L]$.
In practice, the road is discretized into a polyline of $L$ control points $\mathcal{P} = \{\mathbf{p}_1, \dots, \mathbf{p}_L\}$, where $\mathbf{p}_i \in \mathbb{R}^2$.
The objective of test prioritization is to learn a scoring function $f: \mathcal{P} \to [0, 1]$ that estimates the probability of simulation failure, inducing a permutation $\pi$ that sorts test cases in descending order of predicted risk.

The physical vehicle dynamics in an autonomous driving simulation (friction, tire-road lateral forces, and lane deviation) are completely independent of the arbitrary world coordinate frame in which the road is embedded.
Formally, let $g = (\mathbf{R}, \mathbf{t}) \in \mathrm{SE}(2)$ denote a rigid planar motion consisting of a rotation matrix $\mathbf{R} \in \mathrm{SO}(2)$ and a translation vector $\mathbf{t} \in \mathbb{R}^2$.
The transformed road points are given by:
\begin{equation}
T_g(\mathbf{p}_i) = \mathbf{R} \mathbf{p}_i + \mathbf{t}.
\end{equation}
A prioritization scoring model $f$ is defined to be \emph{$\mathrm{SE}(2)$-invariant} if and only if:
\begin{equation}
f(T_g(\mathcal{P})) = f(\mathcal{P}) \quad \forall g \in \mathrm{SE}(2).
\label{eq:se2_def}
\end{equation}
When Eq.~\eqref{eq:se2_def} holds, the resulting ranking $\pi$ is identical under all rotations and translations, precluding artificial ranking variance.

\subsection{Intrinsic Coordinate-Free Feature Set}

Prior deep architectures incorporate absolute headings $\theta_i$ or trigonometric components $(\sin\theta_i, \cos\theta_i)$.
Under a planar rotation $\mathbf{R}_\alpha$ by angle $\alpha$, the heading transforms as $\theta'_i = \theta_i + \alpha$, which directly shifts input activations and destroys invariance.

To achieve exact invariance, SE2RoadNet extracts seven differential geometric invariants that depend solely on intrinsic curve geometry (Table~\ref{tab:features}).
By the fundamental theorem of planar curves~\cite{frenet}, any smooth curve is uniquely determined up to an $\mathrm{SE}(2)$ transformation by its signed curvature $\kappa(s)$.
Given three consecutive points $\mathbf{p}_{i-1}, \mathbf{p}_i, \mathbf{p}_{i+1}$, we compute the signed Menger curvature:
\begin{equation}
\kappa_i = \frac{4 A_i}{a_i b_i c_i} \cdot \mathrm{sgn}\left(\det[\mathbf{p}_i - \mathbf{p}_{i-1}, \;\mathbf{p}_{i+1} - \mathbf{p}_i]\right),
\end{equation}
where $a_i, b_i, c_i$ are triangle side lengths and $A_i$ is the triangle area computed via Heron's formula.
The complete feature vector at control point $i$ is $\mathbf{x}_i = [f_0(i), \dots, f_6(i)]^\top \in \mathbb{R}^7$.
Crucially, every channel $f_k(i)$ satisfies $f_k(T_g(\mathbf{p}_i)) = f_k(\mathbf{p}_i)$ for all $g \in \mathrm{SE}(2)$.

\begin{table}[t]
\caption{The 7-Channel Coordinate-Free Feature Set of SE2RoadNet.}
\label{tab:features}
\centering{\small
\setlength{\tabcolsep}{4pt}
\begin{tabular}{clll}
\toprule
\textbf{Ch.} & \textbf{Feature} & \textbf{Mathematical Definition} & \textbf{Physical Meaning} \\
\midrule
$f_0$ & Segment length & $\ell_i = \|\mathbf{p}_{i+1} - \mathbf{p}_i\|_2$ & Local discretization \\
$f_1$ & Turn angle & $|\Delta\theta_i| = \arccos\left(\frac{\mathbf{v}_i \cdot \mathbf{v}_{i+1}}{\|\mathbf{v}_i\|\|\mathbf{v}_{i+1}\|}\right)$ & Bending intensity \\
$f_2$ & Signed curvature & $\kappa_i = 4 A_i / (a_i b_i c_i) \cdot \mathrm{sgn}(\Delta)$ & Lateral acceleration factor \\
$f_3$ & Curvature jerk & $\Delta\kappa_i / \ell_i = (\kappa_{i+1} - \kappa_i) / \ell_i$ & Rate of steering change \\
$f_4$ & Curvature accel. & $\Delta^2\kappa_i / \ell_i^2$ & Steering impulse \\
$f_5$ & Arclength norm. & $s_i / L_{\text{tot}} = \sum_{j=1}^i \ell_j / L_{\text{tot}}$ & Relative trajectory progress \\
$f_6$ & Local curv. std & $\mathrm{std}(\kappa_{i-w:i+w})$ & Curvature volatility ($w=5$) \\
\bottomrule
\end{tabular}}
\end{table}

\subsection{SE2RoadNet Transformer Architecture}

Fig.~\ref{fig:arch_se2} illustrates the complete pipeline of SE2RoadNet.

\begin{figure*}[t]
\centering
\begin{tikzpicture}[
  font=\sffamily\scriptsize,
  >=Stealth,
  node distance=0.36cm,
  arr/.style={->, thick},
  blk/.style={draw, rounded corners=2pt, fill=#1, align=center,
              inner sep=3.5pt, minimum width=2.9cm, minimum height=0.82cm},
  grp/.style={draw=#1!60, dashed, rounded corners=5pt,
              inner sep=7pt, fill=#1!7},
  htitle/.style={font=\sffamily\bfseries\footnotesize},
]

%% ===== LEFT: INPUT & SE(2)-INVARIANT FEATURE EXTRACTION =====
\node[blk=blue!20, minimum width=3.0cm] (sdc) at (0,0)
    {SDC Test Case\\raw road pts $(x_i, y_i)$};

\node[blk=orange!25, below=0.80cm of sdc] (resamp)
    {Uniform Resampling\\$L{=}197$ points};

\node[blk=orange!25, below=0.36cm of resamp, minimum height=0.95cm] (feat)
    {Extract \textbf{7-Ch Invariant}\\Features $(f_0{-}f_6)$\\
     \textcolor{red!70!black}{\tiny $\times$ no $(x,y)$, no abs.\ heading}};

\node[draw, fill=white, inner sep=3pt, below=0.36cm of feat,
      font=\sffamily\tiny, minimum width=2.85cm] (ftab)
    {\begin{tabular}{@{}ll@{}}
       $f_0$ & Segment length $\ell$ \\
       $f_1$ & $|\Delta\theta|$ (intrinsic angle) \\
       $f_2$ & Signed curvature $\kappa$ \\
       $f_3$ & $d\kappa/ds$ (curvature jerk) \\
       $f_4$ & $d^2\kappa/ds^2$ (curvature accel.) \\
       $f_5$ & $s/L$ (normalized arclength) \\
       $f_6$ & Local std of $\kappa$ \\
     \end{tabular}};

\node[blk=orange!40, below=0.36cm of ftab, minimum width=3.0cm, minimum height=0.9cm] (znorm)
    {$(197{\times}7)$\\z-normalized matrix};

\begin{pgfonlayer}{background}
  \node[grp=orange, fit=(resamp)(feat)(ftab)(znorm),
        label={[htitle]above=2pt:SE(2)-Invariant Feature Extraction}] (leftgrp) {};
\end{pgfonlayer}

%% ===== MIDDLE: SE2ROADNET ARCHITECTURE =====
\node[blk=teal!20, minimum width=3.2cm, text width=3.2cm] (proj) at (6.2,0)
    {Linear + LN + GELU\\$7 \to 192$};

\node[blk=teal!20, minimum width=3.2cm, text width=3.2cm, below=0.36cm of proj] (cls)
    {Append \texttt{[CLS]} token\\(length $L+1 = 198$)};

\node[blk=yellow!35, minimum width=3.3cm, minimum height=1.45cm, below=0.45cm of cls] (enc)
    {\begin{tabular}{c}
      \textbf{Invariant Transformer Block $\times 6$}\\[3pt]
      {\scriptsize LN $\rightarrow$ MHA (8 heads, $d_k{=}24$) + rel-$\Delta s$ bias + res}\\[-2pt]
      {\scriptsize LN $\rightarrow$ FFN (dim 512, GELU) + res}
    \end{tabular}};

\node[blk=purple!20, minimum width=3.2cm, below=0.36cm of enc] (pool)
    {Pool \texttt{[CLS]} representation};

\node[blk=purple!30, minimum width=3.2cm, text width=3.25cm, below=0.36cm of pool, minimum height=0.95cm] (mlp)
    {LN $\to 192 \to 64$\\(GELU, Drop 0.2)\\$\to 1 \to$ Sigmoid};

\node[blk=red!25, minimum width=3.2cm, below=0.36cm of mlp] (yhat)
    {$\hat{y} \in [0,1]$: Failure Probability};

\begin{pgfonlayer}{background}
  \node[grp=teal, fit=(proj)(cls)(enc)(pool)(mlp)(yhat),
        label={[htitle]above:SE2RoadNet Backbone}] (midgrp) {};
\end{pgfonlayer}

%% ===== RIGHT: TRAINING & OBJECTIVES =====
\node[blk=green!25, minimum width=3.2cm, text width=3.1cm] (focal) at (12.5, -0.6)
    {Focal Loss ($\gamma{=}1.5$)\\Class-imbalance robust};

\node[blk=green!25, minimum width=3.2cm, text width=3.1cm, below=0.36cm of focal] (pinn)
    {PINN Auxiliary Loss\\$\mathcal{L}_{\text{PINN}}$ (vehicle dynamics)};

\node[blk=green!25, minimum width=3.2cm, text width=3.1cm, below=0.36cm of pinn] (swa)
    {SWA Averaging\\(epochs 56--80)};

\begin{pgfonlayer}{background}
  \node[grp=green, fit=(focal)(pinn)(swa),
        label={[htitle]above:Physics-Informed Optimization}] (rightgrp) {};
\end{pgfonlayer}

\node[blk=gray!30, minimum width=3.2cm, text width=3.1cm, below=1.0cm of rightgrp.south] (ranker)
    {Sort Descending by $\hat{y}$\\$\to$ Prioritized Execution $\pi$};

%% ===== CONNECTING ARROWS =====
\draw[arr] (sdc) -- (resamp);
\draw[arr] (resamp) -- (feat);
\draw[arr] (feat) -- (ftab);
\draw[arr] (ftab) -- (znorm);

\draw[arr] (znorm.east) -- ++(1.0,0) |- (proj.west);
\draw[arr] (proj) -- (cls);
\draw[arr] (cls) -- (enc);
\draw[arr] (enc) -- (pool);
\draw[arr] (pool) -- (mlp);
\draw[arr] (mlp) -- (yhat);

\draw[arr] (yhat.east) -- ++(1.2,0) |- (rightgrp.west);
\draw[arr] (rightgrp.south) -- (ranker.north);

\end{tikzpicture}
\caption{The end-to-end SE2RoadNet pipeline: Raw road control points are extracted into a 7-channel coordinate-free matrix; a 6-layer Invariant Transformer with relative-arclength attention bias predicts failure probability $\hat{y}$ regularized by a vehicle-dynamics PINN loss and SWA.}
\label{fig:arch_se2}
\end{figure*}

\textbf{Relative-Arclength Attention Bias}:
Standard Transformers utilize absolute sinusoidal or learned positional embeddings $\mathbf{e}_i$, which inject an arbitrary coordinate origin into the sequence.
Instead, SE2RoadNet uses relative position encoding embedded directly into the self-attention logits.
For query token $i$ and key token $j$, the attention weight is computed as:
\begin{equation}
A_{ij} = \frac{\mathbf{q}_i^\top \mathbf{k}_j}{\sqrt{d_k}} + B_{ij},
\end{equation}
where $B_{ij}$ depends strictly on the relative arclength distance $\Delta s_{ij} = |s_i - s_j|$.
To project the scalar distance $\Delta s_{ij}$ into high-dimensional space without imposing artificial polynomial decay, we parameterize $B_{ij}$ using Random Fourier Features (RFF)~\cite{rff}:
\begin{equation}
\mathbf{z}(\Delta s_{ij}) = \sqrt{\frac{2}{D}} \left[ \cos(\omega_1 \Delta s_{ij} + b_1), \dots, \cos(\omega_D \Delta s_{ij} + b_D) \right]^\top,
\end{equation}
where $\omega_d \sim \mathcal{N}(0, \sigma_{\mathrm{rff}}^2)$ and $b_d \sim \mathcal{U}(0, 2\pi)$.
The bias is then obtained via a linear projection $B_{ij} = \mathbf{w}_{\mathrm{bias}}^\top \mathbf{z}(\Delta s_{ij})$.
Because $\Delta s_{ij} = |s_i - s_j|$ is independent of the road's global position, orientation, or starting coordinate, $B_{ij}$ is provably $\mathrm{SE}(2)$-invariant.

\subsection{Physics-Informed Regularization (PINN Loss)}

In realistic driving simulations, vehicle trajectory deviation and out-of-bound (OOB) failures are predominantly triggered by lateral acceleration exceeding tire traction:
\begin{equation}
a_c(s) = v(s)^2 \kappa(s) \le \mu g,
\end{equation}
where $v(s)$ is vehicle velocity, $\mu$ is the tire-road friction coefficient, and $g$ is gravitational acceleration.
When a road features extreme peak curvature $\kappa_{\max} = \max_s |\kappa(s)|$, the required steering angle and lateral force increase monotonically.
A physically valid prioritizer must respect this monotonic relationship: small perturbation increases in $\kappa_{\max}$ should not cause non-monotonic drops in predicted failure probability.

We formulate a soft physics-informed regularizer $\mathcal{L}_{\text{PINN}}$:
\begin{equation}
\mathcal{L}_{\text{PINN}} = \frac{1}{|\mathcal{B}|} \sum_{k=1}^{|\mathcal{B}|} \max\left(0, -\frac{\partial \hat{y}_k}{\partial \kappa_{\max, k}}\right) + \lambda_v \cdot \max\left(0, \hat{y}_k \cdot \mathbb{I}_{\{\kappa_{\max, k} < \kappa_{\text{safe}}\}} - \epsilon\right),
\end{equation}
where the first term penalizes negative gradients of the risk score with respect to peak curvature, and the second term penalizes over-predicting failure on geometrically benign roads where $\kappa_{\max} < \kappa_{\text{safe}}$ ($0.015\,\mathrm{m}^{-1}$).
The overall training objective combines focal classification loss $\mathcal{L}_{\text{Focal}}$ ($\gamma=1.5$)~\cite{focal} with the PINN penalty:
\begin{equation}
\mathcal{L} = \mathcal{L}_{\text{Focal}} + \lambda_{\text{PINN}} \mathcal{L}_{\text{PINN}}.
\end{equation}

\subsection{Stochastic Weight Averaging (SWA)}

To prevent the model from settling into sharp local minima, we employ Stochastic Weight Averaging (SWA)~\cite{swa}.
During training epochs $56$ through $80$, model weights are periodically recorded and averaged:
\begin{equation}
\mathbf{w}_{\mathrm{SWA}} = \frac{1}{K} \sum_{k=1}^K \mathbf{w}_k.
\end{equation}
SWA produces wider, flatter optima that significantly enhance out-of-distribution generalization across diverse simulator road profiles.
